\documentclass[]{standalone}
\usepackage[T1]{fontenc}
\usepackage{tikz}
\usepackage{amsmath, amsfonts}
\usepackage{xintexpr}
\usepackage{xparse}
\usepackage{comment}
\usepackage{ifthen}
\usetikzlibrary{arrows, shapes.geometric, calc, positioning, decorations.pathreplacing, patterns.meta}

\makeatletter
\def\input@path{{./includes/}{./includes/ut/diags/}}
\makeatother

\definecolor{malachite}{rgb}{0.04, 0.85, 0.32}
\definecolor{mikadoyellow}{rgb}{1.0, 0.77, 0.05}
\definecolor{fluorescentpink}{rgb}{1.0, 0.08, 0.58}

\NewDocumentEnvironment{clipThis}{}{
    \begin{scope}
        \clip (-1.2*\offs, -0.1) rectangle (\pgwidth, \cropUpperYRes);
}{
    \end{scope}
}

\newcommand{\drawBlockOutline}[6]{
    \xdef\rawX{#1}
    \xdef\m{#2}
    \xdef\xl{#3}
    \xdef\xr{#4}
    \xdef\yb{#5}
    \xdef\yt{#6}
    \pgfmathsetmacro\rhsStartX{(\xl + \xr)/2}
    \ifnum 0=\m
        \pgfmathsetmacro\rhsStartX{\xl}
    \fi
    \pgfmathtruncatemacro{\drawTop}{mod(\rawX, 2) == 1}

    \begin{clipThis}
        \ifnum 1=\drawTop
            \draw [blockOutline] (\xl, \yt) -- (\xr, \yt);
        \fi
        \draw [blockOutline] (\rhsStartX, \yb) -| (\xr, \yt);
        \draw [blockOutline] (\xl, \yb) -| (\xl, \yt);
    \end{clipThis}
}

\begin{document}%
\begin{tikzpicture}[node distance=0.1cm]%
    \tikzstyle{blockOutline} = [line cap=rect, opacity=1]


    \def\maxMu{7}
    \pgfmathtruncatemacro{\maxMuPOneRes}{\maxMu + 1}
    \xdef\maxMuPOne{\maxMuPOneRes}
    \pgfmathtruncatemacro{\hlSuperMu}{\maxMu - 2}
    \pgfmathtruncatemacro{\nBlocks}{2^(\maxMu+1) - 31}
    \pgfmathtruncatemacro{\proofAnchor}{\nBlocks-6}
    \pgfmathtruncatemacro{\focusBlock}{37}
    \xdef\demoBlock{136}
    \xdef\pgwidth{14}  % cm
    \pgfmathsetmacro\blockW{\pgwidth / \nBlocks}
    \pgfmathsetmacro\offsRes{\blockW/2}
    \xdef\offs{\offsRes}
    \xdef\blockH{.75}
    \pgfmathsetmacro\cropUpperYRes{\maxMuPOne*\blockH + 0.02}

    % * variant: draw top
    % args: \m x1 x2 y1 y2

    \foreach \m in {0,...,\maxMu}{
        \pgfmathtruncatemacro{\nBlocksAt}{ceil(\nBlocks / (2^\m))}
        \pgfmathtruncatemacro{\twoTTMu}{2^\m}
        \pgfmathsetmacro\thisBWidth{\blockW * 2^\m}
        \pgfmathtruncatemacro\nextM{\m + 1}
        \foreach \x in {0,...,\nBlocksAt}{
            \pgfmathtruncatemacro\nextX{\x + \twoTTMu}
            \pgfmathsetmacro\xStart{(\x * \thisBWidth)}
            % \pgfmathsetmacro\xEnd{min(\x * \thisBWidth + \thisBWidth, \pgwidth - \offs)}
            \pgfmathsetmacro\xEnd{(\x * \thisBWidth + \thisBWidth)}
            \pgfmathtruncatemacro\isOkay{\xStart < \xEnd}
            \ifnum 1=\isOkay
                % block numbers to focus on
                \pgfmathtruncatemacro\lower{(2*\x-5)*2^(\m-1)}
                \pgfmathtruncatemacro\upper{(2*\x+5)*2^(\m-1)}
                \pgfmathtruncatemacro\upperHead{(2*\x+8)*2^(\m-1)}

                \pgfmathtruncatemacro\myHeight{\x*2^\m}
                \pgfmathtruncatemacro\nextHeight{(\x+1)*2^\m}
                \pgfmathtruncatemacro\prevSBH{(\x-1)*2^\m}
                \pgfmathtruncatemacro\nextSBH{(\x+1)*2^\m}

                \pgfmathtruncatemacro\isFocusBlock{(\prevSBH < \focusBlock && \myHeight >= \focusBlock)}
                \pgfmathtruncatemacro{\isMainProof}{\proofAnchor <= \upperHead && (\myHeight <= \proofAnchor || (\myHeight > \proofAnchor && \m == 0))}
                \def\blockStyle{}
                % check if this block will be used as part of the NIPoPoW
                \pgfmathtruncatemacro{\isSuper}{0 && mod(\x, 2^(\hlSuperMu-\m+1)) == 0}
                % \pgfmathtruncatemacro{\setBlockStyle}{\isSuper || \isFocusBlock || \isMainProof}

                \pgfmathtruncatemacro{\isDemonstrationBlock}{\m < 4 && \myHeight == \demoBlock}

                % (\m > (\maxMu - 2) && \myHeight > (\proofAnchor-16)) &&
                \pgfmathtruncatemacro{\shouldDrawAtAll}{(\myHeight <= \nBlocks)}


                \ifnum 1=\shouldDrawAtAll
                    \def\mpStyle{}
                    \def\focusStyle{}

                    \ifnum 1=\isMainProof
                        \def\mpStyle{fill=mikadoyellow}
                    \fi
                    \ifnum 1=\isDemonstrationBlock
                        \def\mpStyle{pattern color=gray, pattern={Lines[angle=45, distance=2mm, line width=1mm]}}
                    \fi

                    \ifnum 1=\isFocusBlock
                        \def\focusStyle{pattern color=black, pattern={Hatch[angle=-45, distance=1mm, line width=.3mm]}}

                        \ifnum 0=\isMainProof
                            \path [fill=cyan] (\xStart, \m*\blockH) rectangle (\xEnd, \nextM*\blockH);
                        \else
                            \def\mpStyle{}
                            \begin{clipThis}
                                \path [fill=mikadoyellow, opacity=.4, fill opacity=.75] (\xStart, \m*\blockH) rectangle (\xEnd, \nextM*\blockH);
                            \end{clipThis}
                        \fi
                    \fi
                    \begin{clipThis}
                        % \clip (-\offs, 0) rectangle (\pgwidth, \maxMuPOne*\blockH);
                        \path [opacity=.4, fill opacity=.75, style/.expanded=\mpStyle, style/.expanded=\focusStyle] (\xStart, \m*\blockH) rectangle node (b-\m-\x) {} (\xEnd, \nextM*\blockH);
                    \end{clipThis}
                    \drawBlockOutline{\x}{\m}{\xStart}{\xEnd}{\m*\blockH}{\nextM*\blockH}
                \fi
            \fi
        }
        \pgfmathsetmacro{\muLabelH}{(\m+0.5)*\blockH}
        \node [anchor=east] at (0.1, \muLabelH) {\small $\mu$: \m\;\;\;$\cdots$};
    }

    % label+arrow infix proof target
    \node (tgLabel) [anchor=north east, align=center, xshift=-5mm, yshift=-3mm, font=\small] at (b-0-\focusBlock.south) { Infix Proof \\ Target Block };
    \draw [->, rounded corners=3mm] (tgLabel.east) -| ([yshift=-3mm]b-0-\focusBlock.south);

    % label+arrow suffix proof target
    \node (paLabel) [anchor=north east, align=center, xshift=-5mm, yshift=-3mm, font=\small] at (b-0-\proofAnchor.south) { Suffix Proof \\ Target Block };
    \draw [->, rounded corners=3mm] (paLabel.east) -| ([yshift=-3mm]b-0-\proofAnchor.south);

    % l+a explainer-block
    % , xshift=-5mm
    \node (ebLabel) [anchor=north, align=center, yshift=-3mm, font=\small] at (b-0-\demoBlock.south) { Blocks directly above another \\ block are the same block };
    % loop through outline of blocks
    % track the last top left and top right points of each block as we ascend
    % start from the midpoint of bottom block.
    \xdef\lastTL{(\demoBlock*\blockW, 0)}
    \xdef\lastTR{(\demoBlock*\blockW, 0)}
    \xdef\x{\demoBlock*\blockW}
    % \foreach \m in {0,...,3} {
    %     \pgfmathsetmacro{\thisBWidth}{\blockW * 2^\m}
    %     \pgfmathsetmacro{\leftX}{\x - \thisBWidth / 2}
    %     \pgfmathsetmacro{\rightX}{\x + \thisBWidth / 2}
    %     \pgfmathsetmacro{\nextM}{\m + 1}
    %     % gray shading
    %     \path [fill=fluorescentpink, fill opacity=.25] (\leftX, \nextM*\blockH) rectangle (\rightX, \m*\blockH);
    %     % outline components
    %     \draw [ line cap=rect] \lastTL -- (\leftX, \m*\blockH) -- (\leftX, \nextM*\blockH);
    %     \draw [ line cap=rect] \lastTR -- (\rightX, \m*\blockH) -- (\rightX, \nextM*\blockH);
    %     % for next iter
    %     \xdef\lastTL{(\leftX, \nextM*\blockH)}
    %     \xdef\lastTR{(\rightX, \nextM*\blockH)}
    %     % draw top if we're done
    %     \ifnum 3=\m
    %         \draw[] \lastTL -- \lastTR;
    %     \fi
    % }
    % \draw [->, rounded corners=3mm] (ebLabel.east) -| ([yshift=-3mm]b-0-\demoBlock.south);

    % solid rectangle around suffix proof target
    \node (paTR) [xshift=0.5*\blockW, yshift=0.5*\blockH] at (b-0-\proofAnchor.center) {};
    \node (paBL) [xshift=-0.5*\blockW, yshift=-0.5*\blockH] at (b-0-\proofAnchor.center) {};
    \draw [] (paTR.center) rectangle (paBL.center);
\end{tikzpicture}%
\end{document}
